\section{Repräsentation von Informationen (TiK)}
Ein wichtiger Aspekt an Hochschule ist die Repräsentation von Informationen. Hier spielt sowohl das Erscheinungsbild nach außen als auch die Repräsentation von Informationen innerhalb der einzelnen Bereiche eine essenzielle Rolle. Für diese externe und interne Kommunikation ist die Presse- und Öffentlichkeitsarbeit der Hochschule zuständig (siehe Kapitel \ref{praesidium_label} auf Seite \pageref{praesidium_label}). Nachfolgend wird auf die Teilbereiche   Studentengewinnung, Corporate Identiy und das Handling von Bewerberdaten näher Bezug genommen. 

\subsection{Studentengewinnung}
Wie in Kapitel \ref{praesidium_label} auf Seite \pageref{praesidium_label} bereits erwähnt, obliegt die Zuständigkeit des Marketings dem Präsidialbüro. Generell ist die Studentengewinnung wie folgt aufgeteilt: Zum Einen existiert eine zentrale Studienberatung, bei der sich Interessenten direkt informieren können und zum Anderen werden regelmäßig Besuche von Mitarbeitern der Hochschule an Schulen in der Region durchgeführt. Mitglieder der Fachbereiche berichten vor Ort in den Schulen über die Inhalte der jeweiligen Studiengänge. Neben diesen beiden Maßnahmen zur Studentengewinnung verfügt die Hochschule Emden/Leer ebenso über Onlinemedien, mit deren Hilfe sich die Interessenten über den gewünschten Studiengang und den generellen Ablauf des  Studiums informieren können. 


\subsection{Corporate Identity}
Die Hochschule verfügt über eine CD (Corporate Design) Regelung, welche auf der Webseite öffentlich eingesehen werden kann. Diese wird vom Bereich Marketing zur Verfügung gestellt und gepflegt.\footnote{\url{http://www.hs-emden-leer.de/einrichtungen/praesidialbueropresse-und-oeffentlichkeitsarbeit/corporate-design.html}}

Die CD-Reglung umfasst unter anderem einen CD-Regelungsguide\footnote{\url{http://www.hs-emden-leer.de/fileadmin/user_upload/Einrichtungen/Praesidialbuero/Downloads/HS_CD_Manual_08_2013.pdf}}, sowie diverse Vorlagen für PowerPoint Präsentationen bis hin zu allgemeinen Logos sowie angepasste Logos für jeden Fachbereich (siehe Abbildung \ref{fig_logo_allgemein} und \ref{fig_logo_fb_technik}).

\begin{figure}[h!]
	\centering
	\includegraphics[width=8cm]{kapitel/gruppe2/bilder/hs_logo_allgemein}
	\caption{Allgemeines Logo der Hochschule Emden/Leer}
	\label{fig_logo_allgemein}
\end{figure}

\begin{figure}[h!]
	\centering
	\includegraphics[width=8cm]{kapitel/gruppe2/bilder/hs_logo_technik}
	\caption{Logo des Fachbereiches Technik}
	\label{fig_logo_fb_technik}
\end{figure}

\subsection{Handling von Bewerberdaten}
Durch das Interview mit Herrn Günter Müller konnte festgestellt werden, wie das allgemeine Handling von Bewerberdaten aus Sicht des Rechenzentrums stattfindet. 
Der Bewerber meldet sich mit seinen bei Registrierung erstellten Daten am Hochschulinformationssystem (HIS) der Hochschule an und der generierte Account ist nur für die Dauer des Bewerbungszeitraumes aktiv. Im Fall das der Bewerber nicht angenommen wird, erfolgt im Anschluss die Löschung seines Accounts. Bei Immatrikulation des Bewerbers wird der vorhandene temporäre Account in einen permanenten Account mit erweiterten Informationseingaben, wie z.B. Krankenkassendaten umgewandelt. Auf diese Weise ist sichergestellt, das nur aktive Studenten einen Account zur Verfügung gestellt bekommen. 